% !TeX spellcheck = en_US

\subsection{Motivation}

% \xkcdframe{1172} % version/workflow

\subsection{Configuration Management}
\begin{frame}{Recap -- \insertsubsection{} \mytitlesource{\sommerville}}
	\slideConfigManagement
\end{frame}

\subsection{Version Control}
\begin{frame}<1-4>{Recap -- \insertsubsection{}  \mytitlesource{\sommerville}}
	\slideVersionControl
\end{frame}

\subsection{System Building}
\begin{frame}{\insertsubsection{}  \mytitlesource{\sommerville}}
	\begin{fancycolumns}[animation=none]
		\begin{definition}{System Building}
			Creation of a \textbf{complete executable system} through the \textbf{compilation and linking} of \textbf{system components, external libraries, configuration files}, \dots
		\end{definition}\pause

		\begin{note}{Common Functionalities}
			\begin{itemize}
				\item Automated build script creation
				\item VCS integration
				\item Minimal re-compilation
				\item Build process automation
				\item Test automation
				\item Report creation
				\item Document creation (commonly documentation, \dots)
			\end{itemize}
		\end{note} \pause
		\nextcolumn
		\centering\pic[width=\columnwidth]{configmanagement/system-building}% Diagramm Sommerville 826
	\end{fancycolumns}
\end{frame}

\begin{frame}
	\frametitle{\insertsubsection{} for Java -- Maven}
	\tikzstyle{pomcontent} = [->,dashed,green]
	\tikzstyle{storeaction} = [->,dashed,orange]

	\begin{figure}
		\centering
		\begin{tikzpicture}[node distance=0.5cm, auto]
			\node[draw,ellipse] (maven) {\textbf{Maven}};

			\node[draw, rectangle, above left =1cm and 2cm of maven, minimum width=2.5cm, minimum height=0.7cm] (sourcecode) {Source Code};
			\node[draw, rectangle, below=0.5cm of sourcecode, minimum width=2.5cm, minimum height=0.7cm] (pomxml) {pom.xml};
			\node[draw, rectangle, below=0.5cm of pomxml, minimum width=2.5cm, minimum height=0.7cm] (resources) {Data/Resources};

			\node[draw, rounded corners, align=left, font=\footnotesize\ttfamily, left=0.7cm of pomxml] (dependency) {
				<dependency> \\
				<groupId>junit</groupId> \\
				<artifactId>junit</artifactId> \\
				<version>5.11</version> \\
				</dependency>%
			};%

			\node[draw, rounded corners, align=left, font=\footnotesize\ttfamily, above=0.2cm of dependency] (pomproject) {
				<project> \\
				<modelVersion>4.0.0</modelVersion> \\
				<groupId>org.tubs</groupId> \\
				<artifactId>my-app</artifactId> \\
				<version>0.1.0</version> \\
				</project>%
			};%

			\node[draw, rounded corners, align=left, font=\footnotesize\ttfamily, below=0.2cm of dependency] (pomconfig) {
				<configuration> \\
				<archive> \\
				<manifest> \\
					<mainClass> \\
						org.tubs.Main \\
					</mainClass> \\
				</manifest> \\
				</archive> \\
				</configuration>%
			};%


			\node[draw, dashed, fit=(sourcecode) (pomxml) (resources), inner sep=0.3cm, label=above:\textbf{Project}] (project) {};

			\node[draw, rectangle, above= 1cm of maven, minimum width=2cm, minimum height=1.2cm] (localrepo) {
				MVN Local Repo
			};

			\node[draw, rectangle, below= 1cm of maven, minimum width=2cm, minimum height=1.2cm] (remoterepo) {
				MVN Remote Repo
			};

			\node[draw, rectangle, above right= 1cm and 2cm of maven, minimum width=2.5cm, minimum height=0.7cm] (compile) {Compile};
			\node[draw, rectangle, below=0.5cm of compile, minimum width=2.5cm, minimum height=0.7cm] (test) {Test};
			\node[draw, rectangle, below =0.5cm of test, minimum width=2.5cm, minimum height=0.7cm] (install) {Install};
			\node[draw, rectangle, below =0.5cm of install, minimum width=2.5cm, minimum height=0.7cm] (deploy) {Deploy};

			% Dashed grouping box for lifecycle
			\node[draw, dashed, fit=(compile) (test) (install) (deploy), inner sep=0.3cm, label=above:\textbf{Lifecycle}] (lifecycle) {};

			% Dashed grouping box for repositories

			% Arrows
			\draw[->, thick] (project.east |- maven.west) -- (maven.west);
			\draw[->, thick] (localrepo) -- (maven);
			\draw[->, thick] (remoterepo) -- (maven);
			\draw[->, thick] (maven.east) -- (lifecycle.west |- maven.east);


			% pom arrows
			\draw[pomcontent] (dependency.east) -- (pomxml.west);
			\draw[pomcontent] (pomproject.east) -- (pomxml.north west);
			\draw[pomcontent] (pomconfig.east) -- (pomxml.south west);
			\draw[pomcontent, bend left=30] (pomxml.east) -- (remoterepo.west);

			% action arrows
			\draw[storeaction] (deploy.west) -- (remoterepo.east);
			\draw[storeaction] (install.north west) to[bend right=30] (localrepo.south east);

		\end{tikzpicture}
	\end{figure}
\end{frame}

\subsection{Change Management}
\begin{frame}{\insertsubsection{}  \mytitlesource{\sommerville}}
	\begin{fancycolumns}[animation=none]
		\begin{definition}{Change Management}
			\textbf{Handling, documentation, and prioritization} of \textbf{desired changes} to a software product, as well their \textbf{structured  realization}
		\end{definition}\pause

		\begin{note}{Assessment criteria for change requests}
			\begin{itemize}
				\item Impact of non-acceptance
				\item Advantages of change
				\item Number of impacted users
				\item Costs of change
				\item Impact on product release cycle
			\end{itemize}
		\end{note} \pause
		\nextcolumn
		\begin{note}{Exemplary Change Management Process in Open Source Projects}
			\begin{enumerate}[<+->]
				\item Issue creation by user or developer
				\item Issue replication, categorization, and prioritization
				\item Issue assignment
				\item Implementation in feature branch
				\item Merge request opened
				\item Review of merge request by core developers
				\item Possibly multiple change/review cycles
				\item Merging of changes, closing issue
				\item Documentation for later release notes
			\end{enumerate}
		\end{note}
	\end{fancycolumns}
\end{frame}

\subsection{Release Management}
\begin{frame}{\insertsubsection{}  \mytitlesource{\sommerville}}
	\begin{fancycolumns}[animation=none]
		\begin{definition}{Release Management}
			\textbf{Preparation} of software product for \textbf{external release} and  \textbf{monitoring} of \textbf{already released} software versions
		\end{definition}\pause

		\begin{note}{Typical content of a release}
			\begin{itemize}
				\item Executable system (e.g., code, \dots)
				\item Configuration files
				\item Resources (e.g., data, localization files, \dots)
				\item Installation wizard
				\item System documentation
				\item (sometimes) packaging and advertisement
			\end{itemize}
		\end{note}\pause
		\nextcolumn
		\vspace{-7mm}
		\begin{note}{Factors for release date planning}
			\begin{itemize}
				\item Competition
				\begin{itemize}
					\item Provide same/more features than competitors
					\item Market position dependent on faster release than competitors (e.g., Corona app)
				\end{itemize}
				\item Marketing
				\begin{itemize}
					\item Advertised release
					\item Convince to update
				\end{itemize}
				\item Platform changes
				\begin{itemize}
					\item e.g., OS update
				\end{itemize}
				\item Technical quality of system
				\begin{itemize}
					\item Detected exploits
					\item Critical bugs
				\end{itemize}
			\end{itemize}
		\end{note}
	\end{fancycolumns}
\end{frame}

%Irgendwann Parallel Development
